\documentclass[10pt,utf8]{beamer}

\mode<presentation> {
%  \usetheme{Boadilla}
  \usetheme{Madrid}
%	\usetheme{Fzu}
  \setbeamercovered{transparent}
}

\usepackage{palatino}
\usepackage{graphicx}
\usepackage{array}
\usepackage{color}
\usepackage{subfigure}
\usepackage{colortbl}
\usepackage{amsmath}
\usepackage{hyperref}
\usepackage{listings}
\usepackage{fancyvrb}

%\setbeamertemplate{caption}{\raggedright\insertcaption\par} %turn off caption prefix ("Figure")

\title{Introduction to systemtap}
\author{Vojtěch Juránek}
\institute[Red Hat]{oVirt storage team}
\date{3.~12.~2019}

\lstdefinestyle{Bash}{
	basicstyle          = \large\ttfamily,
	language            = Bash,
	numbers             = left,
	numberstyle         = \small,
	stepnumber          = 1,
	numbersep           = 5pt,
	backgroundcolor     = \color{white},
	showspaces          = false,
	showstringspaces    = false,
	showtabs            = false,
	frame               = single,
	tabsize             = 2,
	captionpos          = b,
	breaklines          = true,
	breakatwhitespace   = false,
	morestring          = [b]",
	stringstyle         = \color{blue},
	keywordstyle        = \color{magenta},
	commentstyle        = \color{gray},
	identifierstyle     = \color{black},
	moredelim           = **[is][\bfseries]{`}{`},
	moredelim           = **[is][\color{magenta}]{!}{!}, 
	fancyvrb            = true,
}


\begin{document}

\begin{frame}
	\titlepage
\end{frame}

\begin{frame}
	\frametitle{What is }
	\begin{figure}
            \centering
            \includegraphics[height=5cm]{./img/stap.eps}
        \end{figure}
	
	\hspace{1cm}
	
	\begin{itemize}
		\item Tool for probing and gathering information about running Linux system.
		\item Also scripting language for writing systemtap probes.
	\end{itemize}

\end{frame}

\begin{frame}
	\frametitle{How it works}
	\begin{itemize}
		\item Probe script is translated into C.
		\item C program is compiled into loadable kernel module.
		\item Loaded into the kernel at the beginning of the script.
		\item Unloaded from the kernel once script finishes.
	\end{itemize}
\end{frame}

\begin{frame}[fragile]
    \frametitle{Setup}
    \begin{itemize}
        \item Install \texttt{systemtap} package: \texttt{dnf install systemtap}
        \item Eventually install requested version of \texttt{kernel-devel} package if systemtap asks for it.
        \item Create "hello world" script  \texttt{hello.stp}:
            \begin{lstlisting}[style=Bash]
probe begin {
  printf("Hello world!\n")
  exit()
}
            \end{lstlisting}
        \item And run it: \texttt{stap hello.stp}
        \item You can use also shebang line in the script and make it executable: \texttt{\#!/usr/bin/env stap}
    \end{itemize}
\end{frame}

\begin{frame}
    \frametitle{stapprobes}
    Probe points - events, which trigger the probe. Many of them, such as:
    \begin{itemize}
        \item \texttt{begin} - start of systemtap session,
        \item \texttt{end} - start of systemtap session,
        \item \texttt{kernel.function("\$FUNC\_NAME")} - entry to the function named \texttt{\$FUN\_NAME},
        \item \texttt{syscall.\$SYS\_CALL} - entry to the system call \texttt{\$SYS\_CALL},
        \item \texttt{syscall.\$SYS\_CALL.return} - return from the system call \texttt{\$SYS\_CALL},
        \item \texttt{time.ms(\$MS)} - triggers every \texttt{\$MS} milliseconds.
    \end{itemize}
    
    \hspace{0.5cm}
    
    There are many of them, features like support of wildcards, etc., see \texttt{man stapprobes} for more.
\end{frame}

\begin{frame}
    \frametitle{Contextual functions}
    Again, many of them, few examples:
    \begin{itemize}
        \item \texttt{pid()} - process ID of current thread,
        \item \texttt{execname()} - name of current process,
        \item \texttt{pp()} - description of probe point being processed,
        \item \texttt{ppfunc()} - function in which the probe was placed,
        \item \texttt{\$\$vars} - list of local variables,
        \item \texttt{print\_backtrace()} - print kernel trace,
        \item \texttt{print\_ubacktrace()} - print user-space trace.
    \end{itemize}
    
    \hspace{0.5cm}
    
    See \texttt{man} pages of  \texttt{apropos function::* } for more details.
\end{frame}

\begin{frame}
    \frametitle{Basics of systemtap scripting language}
    \begin{itemize}
        \item Variables can be just assigned - will be declared and initialized automatically.
        \item Types are inferred automatically.
        \item Variables are by default local.
        \item Global variables can be declared anywhere in the script using keyword \texttt{global}.
        \item Global variable are accessible by all probes.
        \item Target variable - allow to access context information, e.g. \texttt{\$\$vars} ( list of local variables).
    \end{itemize}

\end{frame}
    
\begin{frame}
    \frametitle{Basics of systemtap scripting language}
    \begin{itemize}
        \item \texttt{if (COND) STATEMENT else STATEMENT} - if/else statement.
        \item \texttt{while(COND) STATEMENT} - while loop.
        \item \texttt{for(INIT; COND; INC) STATEMENT} - for loop.
        \item Loops can use \texttt{break} and \texttt{continue}.
    \end{itemize}
\end{frame}

\begin{frame}
    \frametitle{Basics of systemtap scripting language}
    \begin{itemize}
        \item Custom functions (\texttt{function} keyword).
        \item Arrays - must be declared as global.
        \item Support for aggregates, e.g. \texttt{@sum}, \texttt{@avg}, \texttt{@count} etc.
        \item Support for creating custom libraries (tapsets).
        \item Can even embed C code (must be run in "guru" mode).
    \end{itemize}
\end{frame}

\begin{frame}
	\frametitle{Resources}
	\begin{itemize}
		\item \color{blue}\href{http://sourceware.org/systemtap/tutorial.pdf}{systemtap tutorial} \color{black}(presentation is based on this tutorial)
		\item \color{blue}\href{https://sourceware.org/systemtap/SystemTap_Beginners_Guide/}{systemtap beginner guide}
		\item \href{https://sourceware.org/systemtap/examples/}{systemtap examples}
		\item \href{https://sourceware.org/systemtap/langref/}{systemtap language reference}
		\item \href{https://sourceware.org/systemtap/tapsets/}{tapsets}
		\item \href{https://access.redhat.com/documentation/en-us/red\_hat\_enterprise_linux/7/html-single/systemtap\_beginners\_guide/index}{systemtap beginner guide, RHEL 7 documentation}
		\item \href{https://access.redhat.com/documentation/en-us/red_hat_developer\_toolset/8/html/user\_guide/chap-systemtap}{systemtap chapet in RHEL 8 User Guide}
	\end{itemize}
\end{frame}

\begin{frame}
	\centering
	\huge{\textbf{Demo scripts}}
\end{frame}

\begin{frame}[fragile]
    \frametitle{open\_sys.stap}
    \begin{lstlisting}[style=Bash]
probe syscall.open {
  if (substr(filename, 1, 4) == "/sys") {
    printf("%s(%d) open (%s)\n", execname(), pid(), argstr)
  }
}

probe timer.ms(5000){
  exit()
}
    \end{lstlisting}
\end{frame}

\begin{frame}[fragile]
    \frametitle{lvm.stap}
    \begin{lstlisting}[style=Bash]
global pid = -1
probe process("lvm").begin {
  pid = pid()
  printf("Begin: %s pid=%d args=%s\n", execname(), pid, cmdline_args(0, -1, " "))
}

probe process("lvm").end {
  printf("End: %s pid=%d\n", execname(), pid())
}

probe syscall.open {
  if (pid() == pid)
    printf("%s(%d) open (%s)\n", execname(), pid(), argstr)
}
    \end{lstlisting}
\end{frame}

\end{document}
